% Generated draft: requires mathematical and source review.
\documentclass[11pt]{article}
\usepackage[margin=25mm]{geometry}
\usepackage{fontspec}
\setmainfont{Latin Modern Roman}
\usepackage{amsmath,amssymb,mathtools,unicode-math}
\setmathfont{Latin Modern Math}
\usepackage{enumitem,needspace,xurl,hyperref}
\setlength{\parindent}{0pt}
\setlength{\parskip}{5pt}
\setlength{\emergencystretch}{3em}
\newcommand{\N}{\mathbb{N}}
\newcommand{\Z}{\mathbb{Z}}
\newcommand{\Q}{\mathbb{Q}}
\newcommand{\R}{\mathbb{R}}
\newcommand{\C}{\mathbb{C}}
\newcommand{\F}{\mathbb{F}}
\newcommand{\A}{\mathbb{A}}
\newcommand{\PP}{\mathbb{P}}
\newcommand{\E}{\mathbb{E}}
\newcommand{\eps}{\varepsilon}
\newcommand{\dd}{\,\mathrm{d}}
\begin{document}
\section*{UnsolvedMath ID: 20000554}
\textbf{Source problem code:} catalogue entry
\subsection{Definitions and mathematical statement}
The source asks whether partition data for exponential frames depend only on the measure \(|E|\), or also on the geometry of \(E\), and asks for equal-measure sets with different required numbers of pieces or different constants. The ambient measure space, frequency family, type of partition and normalization are not stated in the question itself.

The supplied report uses \(\mathbb T=\R/\Z\) with probability Lebesgue measure, a measurable \(E\subseteq\mathbb T\) of positive measure, and the restricted integer exponentials \(e_n(x)=e^{2\pi i nx}|_E\), \(n\in\Z\), in \(L^2(E)\). A subset \(\Lambda\subseteq\Z\) is a Riesz-sequence index set when there are \(0<a\le b<\infty\) such that
\[
 a\sum_{n\in\Lambda}|c_n|^2\le
 \int_E\left|\sum_{n\in\Lambda}c_ne^{2\pi i nx}\right|^2dx
 \le b\sum_{n\in\Lambda}|c_n|^2
\]
for every finitely supported coefficient family. Let \(r(E)\) be the smallest number of Riesz-sequence pieces in a partition of all \(\Z\), or \(\infty\) if none exists. This definition uses the unnormalized restricted exponentials, each of squared norm \(|E|\).

The report asserts existence of measure-half sets \(E_0,E_*\) such that
\[
 E_0=[0,1/2),\qquad |E_*|=1/2,\qquad r(E_0)=2<r(E_*).
\]
It describes \(E_*\) as an affine normalization of a position-invariant, gap-shifted Kozma--Nitzan--Olevskii recursive construction, and attributes the two-piece obstruction to the Enstad--van Velthoven critical-density theorem and complement density. The recursive set and full theorem hypotheses are not given, so an explicit \(E_*\) cannot be reconstructed here, and its exact partition count above two is not asserted.

For rational grid unions the supplied exact formula is
\[
 E_A=\bigcup_{j\in A}[j/N,(j+1)/N),\quad
 \varnothing\ne A\subseteq\{0,\ldots,N-1\},\qquad
 r(E_A)=\left\lceil\frac N{|A|}\right\rceil.
\]
The report also gives two explicit equal-measure sets
\[
 E_0=[0,1/2),\qquad E_1=[0,1/4)\cup[3/8,5/8),
\]
both with \(r=2\), but different unrestricted optimal common two-piece lower Riesz constants. More precisely, if \(a_E(\Lambda)\) is the infimum of the displayed squared-norm ratio over nonzero finite coefficients, define
\[
 \beta_2(E)=\sup_{\Z=\Lambda_1\sqcup\Lambda_2}
 \min\{a_E(\Lambda_1),a_E(\Lambda_2)\},
\]
with nonempty pieces and no restriction to periodic partitions. The reported conclusion is \(\beta_2(E_0)\ne\beta_2(E_1)\); the two numerical values are absent. It is important that this compares optimal constants over all two-piece partitions, rather than constants from two selected decompositions. The assertions are supplied partial-result claims requiring proof and normalization review.
\subsection{Short English statement}
Can equal-measure sets require different Riesz-sequence partitions of their restricted integer exponential frames? The supplied report says yes, and separately gives explicit equal-measure grid sets with the same partition count but different optimal lower constants. The non-grid construction and numerical constants are not fully supplied.
\subsection{Sources}
\begin{itemize}
\item Local UnsolvedMath registry, permanent ID \(20000554\), source code \texttt{catalogue entry}; AIM analysis corpus, \texttt{aim-analysis-notes.json}, notes[85]. Attribution: Ulam AI / UnsolvedMath Contributors, CC BY 4.0.
\item AIM \emph{Beyond Kadison--Singer: paving and consequences}, source item b13: \url{https://docs.google.com/document/d/1XRw0WyG5PshqboCp99tovCRD1GCK0jXZiNLH_nw618c/edit}. Original remarks and literature fields are empty; no proposer is supplied.
\item The research summary names the Kozma--Nitzan--Olevskii construction and Enstad--van Velthoven critical-density theorem, without full references or statements. It does not give the recursive definition of \(E_*\) or numerical optimal constants.
\item Registry dated August 14, 2026: statement exact, status \texttt{partially\_solved}, classification \texttt{PARTIAL-PROGRESS}; attempt 1 labelled partial result, rigor proved, novelty confidence low. The existence, rational-grid formula and optimal-constant comparison are machine-generated research aids, not peer reviewed, requiring independent verification before citation.
\end{itemize}
\end{document}
